\subsubsection{Linear Maps and Representations}

This subsection blends the two standard viewpoints. Axler emphasises linear maps as the primary objects and matrices as secondary coordinate descriptions; Garcia and Horn emphasises how matrices expose structure. We will use both viewpoints together.

\begin{definition} \hlt{Linear Map}\\
Let $\mathcal{V}$ and $\mathcal{W}$ be vector spaces over the same field $\mathbb{F}$. A function
\begin{align}
T : \mathcal{V} \to \mathcal{W} \nonumber
\end{align}
is linear if for all $\mathbf{u}, \mathbf{v} \in \mathcal{V}$ and all $a,b \in \mathbb{F}$,
\begin{align}
T(a\mathbf{u} + b\mathbf{v}) = aT(\mathbf{u}) + bT(\mathbf{v}). \nonumber
\end{align}
The set of all linear maps from $\mathcal{V}$ to $\mathcal{W}$ is denoted $\LL(\mathcal{V}, \mathcal{W})$. If $\mathcal{V} = \mathcal{W}$, we write $\LL(\mathcal{V})$ and call such a map an operator.
\end{definition}

\begin{definition} \hlt{Graph of a Function}\\
For a function $T:\mathcal{V}\to\mathcal{W}$, the \hlt{graph} of $T$ is the subset of $\mathcal{V}\times\mathcal{W}$ defined by
\begin{align}
\operatorname{graph}(T)\defeq\{(\mathbf{v},T\mathbf{v}):\mathbf{v}\in\mathcal{V}\}. \nonumber
\end{align}
\end{definition}

\begin{proposition} \label{prop:graph-subspace-characterizes-linearity} \hlt{Linearity and Graphs}\\
A function $T:\mathcal{V}\to\mathcal{W}$ is linear if and only if $\operatorname{graph}(T)$ is a subspace of $\mathcal{V}\times\mathcal{W}$.
\begin{proof}
\begin{enumerate}[label=(\roman*)]
\setlength{\itemsep}{0pt}
\item[$(\Rightarrow)$] If $T$ is linear, then $(\mathbf{0},\mathbf{0})\in\operatorname{graph}(T)$. If $(\mathbf{u},T\mathbf{u}),(\mathbf{v},T\mathbf{v})\in\operatorname{graph}(T)$ and $a\in\mathbb{F}$, then
\begin{align}
a(\mathbf{u},T\mathbf{u})+(\mathbf{v},T\mathbf{v})
=(a\mathbf{u}+\mathbf{v},T(a\mathbf{u}+\mathbf{v}))\in\operatorname{graph}(T). \nonumber
\end{align}
Thus $\operatorname{graph}(T)$ is a subspace.
\item[$(\Leftarrow)$] If $\operatorname{graph}(T)$ is a subspace, then for $\mathbf{u},\mathbf{v}\in\mathcal{V}$ and $a\in\mathbb{F}$,
\begin{align}
a(\mathbf{u},T\mathbf{u})+(\mathbf{v},T\mathbf{v})
=(a\mathbf{u}+\mathbf{v},aT\mathbf{u}+T\mathbf{v}) \nonumber
\end{align}
lies in $\operatorname{graph}(T)$. Hence $T(a\mathbf{u}+\mathbf{v})=aT\mathbf{u}+T\mathbf{v}$, so $T$ is linear.
\end{enumerate} \hfill \\
\hlt{Structural Commentary}: A graph records a function as a subset of a product space.\\
\hlt{Key Insight}: Closure of the graph is exactly linearity of the function.
\end{proof}
\end{proposition}

\begin{remark} \hlt{Coordinate-Free Principle}\\
A linear map is completely determined by what it does to a basis. Coordinates do not create the map; they only encode it once bases have been chosen.
\end{remark}

\begin{proposition}
If $(\mathbf{v}_1, \ldots, \mathbf{v}_n)$ is a basis of $\mathcal{V}$ and $\mathbf{w}_1, \ldots, \mathbf{w}_n \in \mathcal{W}$ are arbitrary, then there exists a unique linear map $T : \mathcal{V} \to \mathcal{W}$ such that
\begin{align}
T(\mathbf{v}_j) = \mathbf{w}_j \qquad \text{for } j=1,\ldots,n. \nonumber
\end{align}
\begin{proof}
Every $\mathbf{v} \in \mathcal{V}$ has a unique expression $\mathbf{v} = a_1\mathbf{v}_1 + \cdots + a_n\mathbf{v}_n$. Define
\begin{align}
T(\mathbf{v}) \defeq a_1\mathbf{w}_1 + \cdots + a_n\mathbf{w}_n. \nonumber
\end{align}
This is well-defined because basis coordinates are unique, and linearity is immediate.
\end{proof}
\end{proposition}

\begin{definition} \hlt{Kernel and Range}\\
For $T \in \LL(\mathcal{V}, \mathcal{W})$, define
\begin{align}
\ker(T) \defeq \{ \mathbf{v} \in \mathcal{V} : T(\mathbf{v}) = \mathbf{0} \},
\qquad
\operatorname{range}(T) \defeq \{ T(\mathbf{v}) : \mathbf{v} \in \mathcal{V} \}. \nonumber
\end{align}
\end{definition}

\begin{proposition}
If $T \in \LL(\mathcal{V}, \mathcal{W})$, then $\ker(T)$ is a subspace of $\mathcal{V}$ and $\operatorname{range}(T)$ is a subspace of $\mathcal{W}$.
\end{proposition}

\begin{proposition}
A linear map $T \in \LL(\mathcal{V}, \mathcal{W})$ is injective if and only if
\begin{align}
\ker(T) = \{ \mathbf{0} \}. \nonumber
\end{align}
\end{proposition}

\begin{proposition} \label{prop:direct-sum-map-injective} \hlt{Direct Sum Map}\\
Suppose $\mathcal{U}_1,\ldots,\mathcal{U}_m$ are subspaces of $\mathcal{V}$. Define $\Gamma:\mathcal{U}_1\times\cdots\times\mathcal{U}_m\to\mathcal{U}_1+\cdots+\mathcal{U}_m$ by
\begin{align}
\Gamma(\mathbf{u}_1,\ldots,\mathbf{u}_m)=\mathbf{u}_1+\cdots+\mathbf{u}_m. \nonumber
\end{align}
Then $\mathcal{U}_1+\cdots+\mathcal{U}_m$ is a direct sum if and only if $\Gamma$ is injective.
\begin{proof}
The map $\Gamma$ is linear. Its kernel consists of all $(\mathbf{u}_1,\ldots,\mathbf{u}_m)$ such that $\mathbf{u}_1+\cdots+\mathbf{u}_m=\mathbf{0}$. Thus $\Gamma$ is injective if and only if the only zero decomposition is the trivial one, which is exactly the direct-sum condition.\\
\hlt{Structural Commentary}: The sum map packages all decompositions into one linear map.\\
\hlt{Key Insight}: Directness is injectivity of the map that adds the components.
\end{proof}
\end{proposition}

\begin{theorem} \label{thm:ranknullitythm} \hlt{Rank-Nullity Theorem}\\
If $\mathcal{V}$ is finite-dimensional and $T \in \LL(\mathcal{V}, \mathcal{W})$, then
\begin{align}
\dim(\mathcal{V}) = \dim(\ker(T)) + \dim(\operatorname{range}(T)). \nonumber
\end{align}
\begin{proof}
Let $(\mathbf{u}_1, \ldots, \mathbf{u}_m)$ be a basis of $\ker(T)$ and extend it to a basis
\begin{align}
(\mathbf{u}_1, \ldots, \mathbf{u}_m, \mathbf{w}_1, \ldots, \mathbf{w}_r) \nonumber
\end{align}
of $\mathcal{V}$. We claim that
\begin{align}
(T\mathbf{w}_1, \ldots, T\mathbf{w}_r) \nonumber
\end{align}
is a basis of $\operatorname{range}(T)$. It spans because every vector in $\mathcal{V}$ decomposes into a kernel part plus a combination of the $\mathbf{w}_j$, and it is linearly independent because a relation
\begin{align}
a_1T\mathbf{w}_1 + \cdots + a_rT\mathbf{w}_r = \mathbf{0} \nonumber
\end{align}
would imply $a_1\mathbf{w}_1 + \cdots + a_r\mathbf{w}_r \in \ker(T)$, contradicting linear independence of the basis of $\mathcal{V}$. Therefore
\begin{align}
\dim(\operatorname{range}(T)) = r,
\qquad
\dim(\ker(T)) = m, \nonumber
\end{align}
and $\dim(\mathcal{V}) = m+r$.
\end{proof}
\end{theorem}

\begin{definition} \hlt{Rank and Nullity}\\
For a finite-dimensional domain, the rank of $T$ is $\dim(\operatorname{range}(T))$ and the nullity of $T$ is $\dim(\ker(T))$.
\end{definition}

\begin{definition} \hlt{Invertible Linear Map and Isomorphism}\\
A linear map $T \in \LL(\mathcal{V}, \mathcal{W})$ is invertible if there exists $S \in \LL(\mathcal{W}, \mathcal{V})$ such that
\begin{align}
ST = I_{\mathcal{V}}
\qquad\text{and}\qquad
TS = I_{\mathcal{W}}. \nonumber
\end{align}
An invertible linear map is called an isomorphism, and if one exists we write $\mathcal{V} \cong \mathcal{W}$.
\end{definition}

\begin{proposition}
A linear map is invertible if and only if it is both injective and surjective.
\end{proposition}

\begin{definition} \hlt{Induced Map on a Quotient}\\
For $T\in\LL(\mathcal{V},\mathcal{W})$, define $\widetilde{T}:\mathcal{V}/\ker(T)\to\mathcal{W}$ by
\begin{align}
\widetilde{T}(\mathbf{v}+\ker(T))=T\mathbf{v}. \nonumber
\end{align}
\end{definition}

\begin{theorem} \label{thm:quotient-by-kernel-is-range} \hlt{Quotient by the Kernel}\\
Let $T\in\LL(\mathcal{V},\mathcal{W})$. Then:
\begin{enumerate}[label=(\roman*)]
\setlength{\itemsep}{0pt}
\item $\widetilde{T}$ is a well-defined linear map from $\mathcal{V}/\ker(T)$ to $\mathcal{W}$;
\item $\widetilde{T}$ is injective;
\item $\operatorname{range}(\widetilde{T})=\operatorname{range}(T)$;
\item $\mathcal{V}/\ker(T)\cong\operatorname{range}(T)$.
\end{enumerate}
\begin{proof}
\begin{enumerate}[label=(\roman*)]
\setlength{\itemsep}{0pt}
\item If $\mathbf{v}+\ker(T)=\mathbf{v}'+\ker(T)$, then $\mathbf{v}-\mathbf{v}'\in\ker(T)$, so $T\mathbf{v}=T\mathbf{v}'$. Thus $\widetilde{T}$ is well-defined, and linearity follows from the quotient operations.
\item If $\widetilde{T}(\mathbf{v}+\ker(T))=\mathbf{0}$, then $T\mathbf{v}=\mathbf{0}$, so $\mathbf{v}\in\ker(T)$ and $\mathbf{v}+\ker(T)=\ker(T)$. Hence $\ker(\widetilde{T})=\{\ker(T)\}$.
\item By definition, the values of $\widetilde{T}$ are exactly the values of $T$.
\item Parts (ii) and (iii) show that $\widetilde{T}$ is an isomorphism from $\mathcal{V}/\ker(T)$ onto $\operatorname{range}(T)$.
\end{enumerate} \hfill \\
\hlt{Structural Commentary}: Quotienting by the kernel removes exactly the ambiguity of $T$.\\
\hlt{Key Insight}: Every linear map factors as a quotient followed by an isomorphism onto its range.
\end{proof}
\end{theorem}

\begin{theorem}
If $\mathcal{V}$ is finite-dimensional and $T \in \LL(\mathcal{V})$, then the following are equivalent:
\begin{enumerate}[label=\roman*.]
\setlength{\itemsep}{0pt}
\item $T$ is invertible;
\item $T$ is injective;
\item $T$ is surjective.
\end{enumerate}
\begin{proof}
The implication $(i) \Rightarrow (ii)$ is immediate. If $T$ is injective, then $\ker(T)=\{ \mathbf{0} \}$, so rank-nullity gives
\begin{align}
\dim(\operatorname{range}(T)) = \dim(\mathcal{V}), \nonumber
\end{align}
hence $\operatorname{range}(T)=\mathcal{V}$ and $T$ is surjective. The implication $(iii) \Rightarrow (i)$ follows similarly.
\end{proof}
\end{theorem}

\begin{theorem}
Two finite-dimensional vector spaces are isomorphic if and only if they have the same dimension.
\end{theorem}

\begin{definition} \hlt{Matrix Representation}\\
Let $\beta = (\mathbf{v}_1, \ldots, \mathbf{v}_n)$ be a basis of $\mathcal{V}$ and $\gamma = (\mathbf{w}_1, \ldots, \mathbf{w}_m)$ a basis of $\mathcal{W}$. For $\mathbf{v} \in \mathcal{V}$, denote its coordinate column vector by $[\mathbf{v}]_{\beta}$. For $T \in \LL(\mathcal{V}, \mathcal{W})$, the matrix of $T$ relative to $(\beta,\gamma)$ is the unique matrix $[T]_{\gamma \leftarrow \beta} \in \operatorname{Mat}(m,n,\mathbb{F})$ whose $j$th column is $[T(\mathbf{v}_j)]_{\gamma}$.
\end{definition}

\begin{definition} \hlt{Column Space and Matrix Rank}\\
Let $A=[\mathbf{a}_1\ \cdots\ \mathbf{a}_n]\in\operatorname{Mat}(m,n,\mathbb{F})$. The \hlt{column space} of $A$ is $\operatorname{col}(A)\defeq\operatorname{span}(\mathbf{a}_1,\ldots,\mathbf{a}_n)\subseteq\mathbb{F}^m$, and the \hlt{rank} of $A$ is $\operatorname{rank}(A)\defeq\dim(\operatorname{col}(A))$.
\end{definition}

\begin{proposition}
For every $\mathbf{v} \in \mathcal{V}$,
\begin{align}
[T(\mathbf{v})]_{\gamma} = [T]_{\gamma \leftarrow \beta}[\mathbf{v}]_{\beta}. \nonumber
\end{align}
\begin{proof}
Write $\mathbf{v} = a_1\mathbf{v}_1 + \cdots + a_n\mathbf{v}_n$. Then linearity gives
\begin{align}
T(\mathbf{v}) = a_1T(\mathbf{v}_1) + \cdots + a_nT(\mathbf{v}_n), \nonumber
\end{align}
and the right-hand side is exactly the linear combination of the columns of $[T]_{\gamma \leftarrow \beta}$ with coefficients $a_1,\ldots,a_n$.
\end{proof}
\end{proposition}

\begin{theorem} \label{thm:invertible-matrix-columns-basis} \hlt{Columns of an Invertible Matrix Form a Basis}\\
Let $A=[\mathbf{a}_1\ \cdots\ \mathbf{a}_n]\in\operatorname{Mat}(n,n,\mathbb{F})$ be invertible. Then $(\mathbf{a}_1,\ldots,\mathbf{a}_n)$ is a basis of $\mathbb{F}^n$.
\begin{proof}
The columns span $\mathbb{F}^n$ because for any $\mathbf{y}\in\mathbb{F}^n$, writing $\mathbf{x}=A^{-1}\mathbf{y}$ gives
\begin{align}
\mathbf{y}=A\mathbf{x}=x_1\mathbf{a}_1+\cdots+x_n\mathbf{a}_n. \nonumber
\end{align}
They are linearly independent because if $c_1\mathbf{a}_1+\cdots+c_n\mathbf{a}_n=\mathbf{0}$, then $A\mathbf{c}=\mathbf{0}$, so $\mathbf{c}=A^{-1}\mathbf{0}=\mathbf{0}$. Hence $(\mathbf{a}_1,\ldots,\mathbf{a}_n)$ is a basis of $\mathbb{F}^n$.\\
\hlt{Structural Commentary}: An invertible matrix is a coordinate isomorphism of $\mathbb{F}^n$. Its columns are the images of the standard coordinate directions.\\
\hlt{Key Insight}: Invertibility forces the columns to span and prevents nontrivial column relations.
\end{proof}
\end{theorem}

\begin{lemma} \label{lem:two-sided-inverse-uniqueness} \hlt{Uniqueness of Two-Sided Inverses}\\
Let $A,B,C\in\operatorname{Mat}(n,n,\mathbb{F})$. If $AB=I$ and $BC=I$, then $A=C$.
\begin{proof}
Using $BC=I$ and $AB=I$,
\begin{align}
A=AI=A(BC)=(AB)C=IC=C. \nonumber
\end{align}
\hlt{Structural Commentary}: A common middle inverse forces the left and right inverses to coincide.\\
\hlt{Key Insight}: Associativity turns $AB=I=BC$ into $A=C$.
\end{proof}
\end{lemma}

\begin{theorem} \label{thm:square-one-sided-inverse-is-two-sided} \hlt{One-Sided Inverses of Square Matrices}\\
Let $A,B\in\operatorname{Mat}(n,n,\mathbb{F})$. If $AB=I$, then $BA=I$.
\begin{proof}
Consider the descending sequence of subspaces
\begin{align}
\operatorname{Mat}(n,n,\mathbb{F})\supseteq B\operatorname{Mat}(n,n,\mathbb{F})\supseteq B^2\operatorname{Mat}(n,n,\mathbb{F})\supseteq \cdots . \nonumber
\end{align}
Since $\operatorname{Mat}(n,n,\mathbb{F})$ is finite-dimensional, this sequence stabilizes: for some $k$, $B^k\operatorname{Mat}(n,n,\mathbb{F})=B^{k+1}\operatorname{Mat}(n,n,\mathbb{F})$. Thus $B^k=B^{k+1}C$ for some $C\in\operatorname{Mat}(n,n,\mathbb{F})$.\\
From $AB=I$, induction gives $A^rB^r=I$ for all $r\geq 1$. Hence
\begin{align}
I=A^kB^k=A^kB^{k+1}C=(A^kB^k)BC=BC. \nonumber
\end{align}
By Lemma \ref{lem:two-sided-inverse-uniqueness}, $A=C$, so $BA=BC=I$.\\
\hlt{Structural Commentary}: In finite dimensions, repeatedly multiplying by $B$ cannot produce an endlessly shrinking chain of column spaces. Stabilization produces the missing inverse.\\
\hlt{Key Insight}: A one-sided inverse of a square matrix is automatically two-sided.
\end{proof}
\end{theorem}

\begin{proposition} \hlt{Composition and Matrix Multiplication}\\
If $T \in \LL(\mathcal{U}, \mathcal{V})$ and $S \in \LL(\mathcal{V}, \mathcal{W})$, then
\begin{align}
[ST]_{\delta \leftarrow \alpha} = [S]_{\delta \leftarrow \beta}[T]_{\beta \leftarrow \alpha} \nonumber
\end{align}
for any bases $\alpha$ of $\mathcal{U}$, $\beta$ of $\mathcal{V}$, and $\delta$ of $\mathcal{W}$.
\end{proposition}

\begin{theorem} \hlt{Linear Maps and Matrices Are the Same Data}\\
If $\dim(\mathcal{V}) = n$ and $\dim(\mathcal{W}) = m$, then choosing bases of $\mathcal{V}$ and $\mathcal{W}$ gives a vector-space isomorphism
\begin{align}
\LL(\mathcal{V}, \mathcal{W}) \cong \operatorname{Mat}(m,n,\mathbb{F}). \nonumber
\end{align}
Consequently,
\begin{align}
\dim(\LL(\mathcal{V}, \mathcal{W})) = mn. \nonumber
\end{align}
\end{theorem}

\begin{remark} \hlt{Why This Matters}\\
This is the key bridge between abstraction and computation. The linear map is basis-free; the matrix is what the map looks like after coordinates have been chosen.
\end{remark}

\begin{definition} \hlt{Change of Basis Matrix}\\
Let $\beta$ and $\gamma$ be two bases of the same vector space $\mathcal{V}$. The change-of-basis matrix from $\beta$-coordinates to $\gamma$-coordinates is
\begin{align}
C_{\gamma \leftarrow \beta} \defeq [I_{\mathcal{V}}]_{\gamma \leftarrow \beta}. \nonumber
\end{align}
Thus
\begin{align}
[\mathbf{v}]_{\gamma} = C_{\gamma \leftarrow \beta}[\mathbf{v}]_{\beta} \nonumber
\end{align}
for every $\mathbf{v} \in \mathcal{V}$.
\end{definition}

\begin{proposition}
The change-of-basis matrix $C_{\gamma \leftarrow \beta}$ is invertible, with inverse
\begin{align}
C_{\beta \leftarrow \gamma}. \nonumber
\end{align}
\end{proposition}

\begin{theorem} \hlt{Change of Basis Formula}\\
Let $T \in \LL(\mathcal{V})$, and let $\beta,\gamma$ be two bases of $\mathcal{V}$. Then
\begin{align}
[T]_{\gamma \leftarrow \gamma}
=
C_{\gamma \leftarrow \beta}[T]_{\beta \leftarrow \beta}C_{\beta \leftarrow \gamma}. \nonumber
\end{align}
\begin{proof}
Use the identity
\begin{align}
[T(\mathbf{v})]_{\gamma}
=
C_{\gamma \leftarrow \beta}[T(\mathbf{v})]_{\beta}
=
C_{\gamma \leftarrow \beta}[T]_{\beta \leftarrow \beta}[\mathbf{v}]_{\beta}
=
C_{\gamma \leftarrow \beta}[T]_{\beta \leftarrow \beta}C_{\beta \leftarrow \gamma}[\mathbf{v}]_{\gamma}. \nonumber
\end{align}
for all $\mathbf{v} \in \mathcal{V}$.
\end{proof}
\end{theorem}

\begin{definition} \hlt{Similarity}\\
Matrices $A,B \in \operatorname{Mat}(n,n,\mathbb{F})$ are similar if there exists an invertible matrix $P$ such that
\begin{align}
B = PAP^{-1}. \nonumber
\end{align}
\end{definition}

\begin{remark} \hlt{Similarity as Same Operator, New Coordinates}\\
Similarity does not mean two operators are merely related; it means they are the same operator written in two different bases. This is the structural content behind change of basis.
\end{remark}

\begin{definition} \hlt{Block Matrices}\\
Suppose
\begin{align}
\mathcal{V} = \mathcal{U}_1 \oplus \mathcal{U}_2
\qquad\text{and}\qquad
\mathcal{W} = \mathcal{Z}_1 \oplus \mathcal{Z}_2. \nonumber
\end{align}
If $T \in \LL(\mathcal{V}, \mathcal{W})$, then relative to bases adapted to these decompositions, the matrix of $T$ has block form
\begin{align}
[T] =
\begin{bmatrix}
A_{11} & A_{12} \\
A_{21} & A_{22}
\end{bmatrix}, \nonumber
\end{align}
where each block records the part of $T$ mapping one summand of the domain into one summand of the codomain.
\end{definition}

\begin{proposition} \hlt{Invariant Subspaces and Block Form}\\
Let $T \in \LL(\mathcal{V})$ and suppose
\begin{align}
\mathcal{V} = \mathcal{U} \oplus \mathcal{W}. \nonumber
\end{align}
\begin{enumerate}[label=\roman*.]
\setlength{\itemsep}{0pt}
\item If $\mathcal{U}$ is invariant under $T$, then in a basis adapted to $\mathcal{U} \oplus \mathcal{W}$ the matrix of $T$ has the upper block-triangular form
\begin{align}
\begin{bmatrix}
A & B \\
0 & D
\end{bmatrix}. \nonumber
\end{align}
\item If both $\mathcal{U}$ and $\mathcal{W}$ are invariant under $T$, then the matrix of $T$ has the block-diagonal form
\begin{align}
\begin{bmatrix}
A & 0 \\
0 & D
\end{bmatrix}. \nonumber
\end{align}
\end{enumerate}
\begin{proof}
If $\mathcal{U}$ is invariant, then $T(\mathcal{U}) \subseteq \mathcal{U}$, so no basis vector from $\mathcal{U}$ can acquire a component in $\mathcal{W}$. That forces the lower-left block to be zero. If $\mathcal{W}$ is also invariant, the same argument forces the upper-right block to be zero for the complementary direction as well.
\end{proof}
\end{proposition}

\begin{remark} \hlt{Matrix Viewpoint on Structure}\\
Basic block matrices already capture important geometry. Upper block-triangular form records the presence of an invariant subspace; block diagonal form records a genuine decomposition into independent invariant pieces.
\end{remark}
